\section{Polynomial families and finite exponential mixtures}
\label{sec:algebraic}

The geometric proof works because the dependence of all selected terms
on the parameters has controlled complexity. For a fixed-dimensional
polynomial family, a classical sign-condition theorem provides that
control even when the degrees grow with the original sequence index.
We state the general normalized result first. No semialgebraicity of
the parameter set is needed for this existence theorem.

\begin{theorem}[Uniform hitting for polynomial families]\label{thm:poly-family}
Let $\Theta\subset\R^d$ be nonempty and compact, and let
$f_n\in\R[X_1,\ldots,X_d]$, $n\ge0$, satisfy $f_0=1$ on $\Theta$.
Suppose there are constants $0<\alpha<\beta<1$, $A_0\ge1$, and
an integer $D_0\ge1$ such that, for all $n\ge0$ and $z\in\Theta$,
\begin{equation}\label{eq:family-assumptions}
 \alpha f_n(z)\le f_{n+1}(z)\le\beta f_n(z),
 \qquad \deg f_n\le A_0(n+1)^{D_0}.
\end{equation}
For each $0<p<1$ there exists an open $1$-periodic $H\subset\R$
with $\rho(H)\le6p$ such that
\begin{equation}\label{eq:family-hit}
 \forall x\in\R\ \forall t\in[1,2]\ \forall z\in\Theta
 \quad \exists n\ge1:\ x+tf_n(z)\in H.
\end{equation}
\end{theorem}

By induction, $\alpha^n\le f_n(z)\le\beta^n$, so every member
is strictly positive, strictly decreasing, and converges uniformly to
zero. The family may contain coinciding members and redundant
parameters.

\subsection{The sign-condition theorem used}

\begin{lemma}[Classical full sign-condition bound]\label{lem:signs}
If $S$ real polynomials in $v\ge1$ variables have degrees at most
$\Delta\ge1$, then the number of realized sign vectors in
$\{-1,0,1\}^S$ is at most
\begin{equation}\label{eq:sign-bound}
 \bigl(8\Delta(S+1)\bigr)^v.
\end{equation}
The same bound holds after restricting the variables to an arbitrary
subset of $\R^v$.
\end{lemma}

\begin{proof}[Justification and source]
We use the finite all-sign-condition component bound recorded in
\cite[Section~3.2, equation~(3.3)]{BPR}, with ambient dimension $v$:
\[
 \#\{\text{realized sign vectors}\}
 \le \sum_{j=0}^v\binom Sj4^j\Delta(2\Delta-1)^{v-1}.
\]
A sign realization has at least one connected component, so that
component estimate implies the displayed inequality. Since
$\binom Sj\le S^j$ and
$\sum_{j=0}^v(4S)^j\le(1+4S)^v$, the right-hand side is at most
$2^{v-1}\Delta^v(1+4S)^v\le(8\Delta(S+1))^v$.
Restriction cannot increase the number of realized sign vectors.
This is a finite degree-dependent bound, not just a fixed-degree
asymptotic. It includes zero signs.
\end{proof}

We do not need to distinguish connected components with the same sign
vector: the vector itself determines all the discrete choices below.
This observation also allows an arbitrary compact $\Theta$ in
\cref{thm:poly-family}.

\subsection{Simultaneously encode selection and grid addresses}

Choose $Q,T,c$ as in \eqref{eq:Q-T} and define, for $z\in\Theta$,
\[
 \nu_j(z)=\min\{n\ge1:f_n(z)\le Q^j\},
 \qquad a_j(z)=f_{\nu_j(z)}(z).
\]
Exactly as in \cref{lem:envelopes},
\begin{equation}\label{eq:family-envelope}
 \alpha Q^j<a_j(z)\le Q^j,\quad \nu_j(z)\le Tj,
 \quad a_j(z)-a_{j'}(z)>(\alpha-Q)Q^j\quad(j'>j).
\end{equation}
Thus the common grids, windows, stable centers, and finite random
routing from \cref{sec:synchronization,sec:routing} apply without
change. In particular the fixed-$(z,t)$ failure probability is
$(1-p/2)^{m_0r}$.

Fix $x$ and a center-exposure atom with first default $U$.
There are $m_0r$ relevant pairs $(i,j)$. Put $N=TB$.
For every relevant $j$ and $1\le n\le N$, include the polynomial
\begin{equation}\label{eq:selection-polynomial}
 f_n(z)-Q^j.
\end{equation}
The signs of these polynomials, including equality, determine the
first $n$ at which the level is crossed, and hence determine
$\nu_j(z)$.

For every pair $(i,j)$, each grid boundary
\[
 \gamma\in\Gamma_{b_{e_i}^*}\cap[x,x+2Q^j],
\]
and every $1\le n\le N$, also include
\begin{equation}\label{eq:key-polynomial}
 t f_n(z)-(\gamma-x).
\end{equation}
There are at most $1+4cQ^{-2r}$ such boundaries per pair by
\eqref{eq:boundary-count}. Some of the resulting tests are irrelevant
when $n$ is not selected; retaining them only enlarges the bound.
Once \eqref{eq:selection-polynomial} specifies $\nu_j(z)$, the signs
in \eqref{eq:key-polynomial} for that $n$ determine the selected
point's key. All selected points lie in the interval used to list
the boundaries, by \eqref{eq:family-envelope}.

The combined list contains at most
\begin{equation}\label{eq:family-S}
 S\le m_0rTB\bigl(2+4cQ^{-2r}\bigr)
\end{equation}
polynomials in $v=d+1$ variables $(z,t)$, of degrees at most
\begin{equation}\label{eq:family-degree}
 \Delta\le A_0(TB+1)^{D_0}+1.
\end{equation}
A realized sign vector determines the full local key vector for every
unrevealed table outcome. Representatives can therefore be chosen
before those entries are revealed, exactly as in the geometric proof.
By \cref{lem:signs}, \eqref{eq:family-S}, and
\eqref{eq:family-degree}, their number is at most
\begin{equation}\label{eq:family-complexity}
 C_0(B+1)^{(D_0+2)v}Q^{-2vr},
\end{equation}
where $C_0$ depends on $d,A_0,D_0,T,c,M$, but not on $r_0,r,B$.
Here we used $r\le B$ and absorbed only polynomial factors into
$(B+1)^{(D_0+2)v}$.

\subsection{Complete the normalized family proof}

Choose $M$ so that
\begin{equation}\label{eq:family-theta}
 \theta_v=Q^{-2v}(1-p/2)^{M-1}<1.
\end{equation}
Then choose $d_{\rm tree}$ (the tree height, distinct from parameter
dimension $d$), followed by $g$, so that the no-default probability
and the unstable-center density are each less than $p$.
With these choices fixed, $B(r_0)$ is affine in $r_0$.
Choose $r_0$ large enough that
\[
 C_0(B(r_0)+1)^{(D_0+2)v}\theta_v^{r_0}<p.
\]
The union bound over the representatives now gives the same pointwise
stable-center estimate as \eqref{eq:pointwise-miss}, uniformly for
$(z,t)\in\Theta\times[1,2]$.

Enlarge the random set to $A^+$ and define the residual centers by
\[
 R=\{x:\exists(z,t)\in\Theta\times[1,2],\ 
       x+tf_n(z)\notin A^+\text{ for every }1\le n\le TB\}.
\]
This is closed and periodic because the original polynomials are
continuous and $\Theta\times[1,2]$ is compact. Its expected density
is at most $3p$. The identical open-neighborhood repair gives an open
periodic $H$ of density at most $6p$. On a residual center the original
terms converge to the center, uniformly because $f_n\le\beta^n$.
This proves \cref{thm:poly-family}.

\subsection{Proof for all positive mixtures}

Fix a component count $d\ge1$ and an interval $I_j=[\alpha_j,\beta_j]$
from \eqref{eq:ratio-exhaustion}. Take
\begin{equation}\label{eq:mixture-family}
 \Theta_{d,j}=\{w\in[0,1]^d:\textstyle\sum_iw_i=1\}\times I_j^d,
 \qquad f_n(w,q)=\sum_{i=1}^d w_iq_i^n.
\end{equation}
This is a compact rational semialgebraic parameter set. The polynomials
have degree at most $n+1$, $f_0=1$, and
\[
 \alpha_j f_n(w,q)\le f_{n+1}(w,q)\le\beta_j f_n(w,q).
\]
Thus \cref{thm:poly-family} applies. There is no loss at repeated bases
or zero weights, since all arguments include boundary sign conditions.

Enumerate the pairs $(d,j)$ by $\ell\ge1$ in a computable order.
For each $\ell,k$, with $k\ge0$, choose the normalized blocker for
that family with budget
\[
 p_{\ell,k}=\frac\eps{24}2^{-\ell}2^{-(k+1)}.
\]
Form $C^*$ by the same union of signed contractions as
\eqref{eq:global-C}, and put $F^*_\eps=\R\setminus C^*$.
It is closed, symmetric, periodic, and has density at least $1-\eps/2$.

It remains to verify that a tail can always be brought back into one
of the normalized families. Absorb $\sum_i c_i$ into the dilation,
so assume $\sum_i c_i=1$ and write $a_n=\sum_i c_iq_i^n$.
For $h\ge0$ set
\begin{equation}\label{eq:mixture-tail}
 w_i^{(h)}=\frac{c_iq_i^h}{a_h},
 \qquad
 a_{h+n}=a_h\sum_i w_i^{(h)}q_i^n.
\end{equation}
The new weights remain in the same simplex. Choose $h$ large enough
that $|s|a_h\le1$, then normalize it by a dyadic contraction as in
\eqref{eq:tail-normalization}. The appropriate normalized blocker
gives a hit at an original index $h+n$. Since $h$ can be required to
exceed any prescribed threshold, the hits occur infinitely often.
This proves \cref{thm:mixtures}.

\subsection{An additional class: positive polynomial factors}

\begin{corollary}[Positive exponential polynomials]\label{cor:exp-polynomial}
One common closed, symmetric, $1$-periodic set of measure greater than
$1-\eps$ per unit interval can avoid, with infinitely many omissions,
every nontrivial affine pattern whose terms are
\begin{equation}\label{eq:positive-exp-poly}
 a_n=\sum_{i=1}^d P_i(n)q_i^n,
 \qquad 0<q_i<1,
\end{equation}
where $d$ is finite and the $P_i$ have nonnegative real coefficients,
not all zero. The compact section can also have the effective properties
of \cref{thm:effective-intro}.
\end{corollary}

\begin{proof}
Fix $d$, a degree bound $D$, a base interval $[\alpha,\beta]=I_j\subset(0,1)$,
and an integer $h\ge1$ so large that
$\beta(1+1/h)^D<1$. Consider normalized families
\[
 f_n(w,q)=\sum_{i=1}^d\sum_{k=0}^D
    w_{ik}(1+n/h)^kq_i^n,
 \quad w_{ik}\ge0,\quad\sum_{i,k}w_{ik}=1,
 \quad q_i\in[\alpha,\beta].
\]
They have $f_0=1$ and polynomial parameter degree at most $n+1$.
Each summand's ratio from $n$ to $n+1$ lies between $\alpha$ and
$\beta(1+1/h)^D<1$, and so does the ratio of the sums. Apply
\cref{thm:poly-family}, and assemble blockers over all the countably
many choices of $d,D,j,h$, with summable budgets.

Write $P_i(n)=\sum_k c_{ik}n^k$, $c_{ik}\ge0$. A tail has
\[
 a_{h+n}=A_h\sum_{i,k}w_{ik}^{(h)}(1+n/h)^kq_i^n,
 \quad
 A_h=\sum_{i,k}c_{ik}h^kq_i^h,
 \quad w_{ik}^{(h)}=\frac{c_{ik}h^kq_i^h}{A_h}.
\]
For all $h\ge1$, $A_h>0$; it tends to zero as $h\to\infty$.
Choose $h$ satisfying the contraction condition and $|s|A_h\le1$,
and as large as desired. The normalized family and a dyadic contraction
then produce an arbitrarily late hit. The families have effective
rational polynomial descriptions; the argument of
\cref{sec:effective} therefore applies.
\end{proof}

In \cref{cor:exp-polynomial}, negative polynomial coefficients and
complex bases are outside that corollary. The next theorem removes the
coefficient restriction by proving a separate uniform tail reduction.

\subsection{Signed coefficients and stable positive-root recurrences}

The positivity restriction can in fact be removed when every base is a
positive real number. The price is a countable compact exhaustion that
controls the dominant term before the normalized theorem is applied.

\begin{theorem}[All real exponential polynomials with positive bases]
\label{thm:signed-exp-poly}
For every $0<\eps<1$ there exists one closed, symmetric, $1$-periodic
set $F^{\rm ep}_\eps\subset\R$, of measure at least $1-\eps/2$ in
every unit interval, with the following property. For every finite sum
\begin{equation}\label{eq:signed-exp-poly}
 A_n=\sum_{i=1}^m P_i(n)q_i^n,\qquad q_i\in(0,1),
\end{equation}
with arbitrary real polynomial coefficients, not identically zero
after equal bases have been combined, and every $x\in\R$, $s\ne0$,
infinitely many $x+sA_n$ lie outside $F^{\rm ep}_\eps$.
For rational accuracy its compact section can be effectively closed
with computable measure.
\end{theorem}

\begin{proof}
Combine equal bases and discard zero polynomials. Relabel so that
$q_m$ is the largest remaining base. If $P_m$ has degree $j$ and
nonzero leading coefficient $\lambda$, divide all polynomials by
$\lambda$ and replace $s$ by $s\lambda$. Hence we may suppose
$P_m$ is monic of degree $j$.

Fix integers $m\ge1$, $k\ge j\ge0$, $R\ge1$, and rational
$0<a<b<1$, $0<\delta<b$. Consider the compact parameter set where
all bases belong to $[a,b]$, all polynomial degrees are at most $k$,
all coefficients have absolute value at most $R$, $P_m$ is monic of
degree $j$, and
\[
 q_i\le q_m-\delta\quad(i<m).
\]
For $m=1$ the separation conditions are vacuous. The set is rational
semialgebraic. Empty sets may be skipped by quantifier elimination.
Countably many such boxes cover every normalized nonzero expression.

Put $r_*=1-\delta/b\in(0,1)$. Uniformly on each box, for $n\ge1$,
\begin{equation}\label{eq:dominance}
 A_n=n^jq_m^n(1+E_n),\qquad
 |E_n|\le \frac{Rj}{n}+(m-1)R(k+1)n^k r_*^n.
\end{equation}
Indeed the lower terms of the monic dominant polynomial contribute
at most $Rj/n$ after division by $n^j$. Each other polynomial has
absolute value at most $R(k+1)n^k$, and its base ratio is at most
$1-\delta/q_m\le1-\delta/b=r_*$. Dropping the additional factor
$n^{-j}\le1$ gives the stated estimate.

Choose rational $e\in(0,1)$ such that $b(1+e)/(1-e)<1$.
Then choose $N$ large enough that the right-hand side of
\eqref{eq:dominance} is at most $e$ for every $n\ge N$ and
\[
 \beta_*:=b(1+1/N)^j\frac{1+e}{1-e}<1.
\]
These choices are effective: after imposing
$(1+1/N)^kr_*\le r'_*<1$ for some rational $r'_*\in(r_*,1)$,
the sequence $n^kr_*^n$ is decreasing for $n\ge N$, so the
required tail bounds reduce to rational inequalities at $N$.
It follows that $A_n>0$ for $n\ge N$, and
\begin{equation}\label{eq:dominant-ratios}
 \alpha_*:=a\frac{1-e}{1+e}
 \le\frac{A_{n+1}}{A_n}\le\beta_*<1
 \qquad(n\ge N).
\end{equation}

For \emph{each} integer $h\ge N$, augment this compact box by the
parameter $u$ satisfying $uA_h=1$, and use the sequence of polynomials
\[
 f_n=uA_{h+n},\qquad n\ge0.
\]
The augmented set is compact: $A_h\ge(1-e)h^j a^h>0$ uniformly,
so $u=1/A_h$ is continuous and bounded there. It is rational
semialgebraic, $f_0=1$, and \eqref{eq:dominant-ratios} gives the
required uniform contraction bounds. Parameter degree is at most
$h+n+2$, so \cref{thm:poly-family} applies.

Enumerate all these boxes \emph{and all integers $h\ge N$} as the
normalized families in the signed-contraction assembly of
\cref{sec:global}. The same budgets give one periodic set of density
at least $1-\eps/2$. For any original expression, choose its box and
then $h$ as large as desired with $|s|A_h\le1$. Since
$A_{h+n}=A_h f_n$, a dyadic contraction of the normalized blocker
hits the original affine sequence at an index beyond $h$.
This proves infinitely many omissions. Including every $h\ge N$
is what makes the contractions-only assembly valid for these
normalized tail families.

All enumerated families and tail bounds have effective rational data.
Consequently \cref{prop:effective-H} and the measure modulus in
\cref{sec:effective} prove the last assertion.
\end{proof}

\begin{corollary}[Real linear recurrences with roots in $(0,1)$]
\label{cor:recurrences}
The set $F^{\rm ep}_\eps$ simultaneously avoids, with infinitely many
omissions in every nonzero affine scale, every nonzero real solution of a
homogeneous constant-coefficient linear recurrence whose characteristic
roots are all real and lie in $(0,1)$, allowing repeated roots.
\end{corollary}

\begin{proof}
For completeness, if the distinct roots are $q_i$ with multiplicities
$m_i$, the sequences $n^kq_i^n$, $0\le k<m_i$, satisfy the recurrence.
Applying the shift operator minus $q_i$ lowers the polynomial degree
in the corresponding term. These $\sum_i m_i$ sequences are linearly
independent: in a vanishing combination, divide by the largest base
to isolate its polynomial, which cannot tend to zero unless it is
zero; then repeat for the smaller bases. A recurrence of order
$\sum_i m_i$ has that dimension of solutions, since its initial
values determine it. Therefore every solution has the form
\eqref{eq:signed-exp-poly}. A nonzero solution gives a nonzero
expression, and \cref{thm:signed-exp-poly} applies.
\end{proof}
